% =====================================================================================
% REVISIÓN 1  -  Antecedentes (del problema, de la empresa y teóricos) y tabla comparativa
% Escribe tu texto FUERA de los recuadros \begin{guia} ... \end{guia}.
% =====================================================================================

\section{Antecedentes} \label{sec:antecedentes}

\begin{guia}[Guía: Antecedentes (Revisión 1)]
Este capítulo sitúa tu proyecto: qué se ha hecho antes, en qué organización ocurre y qué conceptos lo sustentan.
Tiene tres apartados (\emph{del problema}, \emph{de la empresa}, \emph{teóricos}) y una \emph{tabla comparativa}.
Escribe aquí, si quieres, un párrafo breve que presente los apartados.
\end{guia}

Este capítulo sitúa el proyecto en tres planos: los trabajos y enfoques existentes sobre el problema, la organización donde se realiza la estadía y el proceso que se busca mejorar, y los conceptos teóricos esenciales que sustentan la propuesta. La comparación objetiva entre los enfoques existentes y el sistema propuesto se resume en la Tabla~\ref{tab:comparativa_antecedentes}.

\subsection{Antecedentes del problema} \label{sec:antecedentes_problema} % audit:req=antecedentes

\begin{guia}[Guía: Antecedentes del problema (Revisión 1)]
\textbf{¿Para qué sirve?} Muestra que conoces lo que ya existe sobre el problema y justifica que haga falta tu trabajo.

\textbf{Debe responder:}
\begin{itemize}[nosep,leftmargin=*]
  \item ¿Cómo ha evolucionado o se ha enfrentado este problema? (breve recorrido histórico o de contexto).
  \item ¿Qué soluciones, sistemas, tesis o artículos existentes abordan el mismo problema o uno parecido (mínimo 3, con
        cita \verb|\cite{clave}|)?
  \item ¿Qué hace cada uno, qué ventajas y qué limitaciones tiene?
  \item ¿Qué queda sin resolver y cómo se conecta con tu propuesta?
\end{itemize}

\textbf{Extensión mínima:} 120 palabras. \textbf{Fuentes:} libros, artículos, documentación técnica oficial o sistemas
reales; evita depender solo de blogs.

\textbf{Errores frecuentes:} listar herramientas sin analizarlas; no citar; copiar descripciones comerciales.

\textbf{Cómo se revisa:} pertinencia de las fuentes, análisis crítico y conexión explícita con el problema.
\end{guia}

% >>> ESCRIBE AQUÍ los antecedentes del problema <<<
Existen múltiples publicaciones y productos donde se documentan diferentes métodos para la digitalización de facturas y documentos comerciales. Cada uno hace uso de distintas técnicas de reconocimiento óptico de caracteres, extracción estructurada y validación fiscal. A continuación se describen tres enfoques representativos, cuyas características se comparan con el sistema propuesto en la Tabla~\ref{tab:comparativa_antecedentes}.

En primer lugar, el enfoque clásico basado en Tesseract combinado con reglas heurísticas constituye la vía abierta más documentada. Diversos trabajos reportan pipelines de preprocesamiento con OpenCV, reconocimiento con Tesseract y extracción mediante expresiones regulares para facturas con formatos conocidos\cite{tesseract_ocr}\cite{opencv_docs}. Su principal fortaleza es la ejecución local sin costos por documento; su principal limitación es la fragilidad de las reglas ante formatos no previstos y la ausencia de validación fiscal específica para el contexto mexicano.

En segundo lugar, los servicios cognitivos en la nube, con Google Cloud Document AI como nube de referencia, proporcionan OCR de alta calidad y extracción estructurada pre-entrenada mediante API\cite{gcp_document_ai}. Su fortaleza radica en la exactitud y en la extracción de tablas; sus limitaciones frente a lo propuesto son la dependencia de un proveedor externo, el costo por documento, el envío de documentos fiscales fuera de la infraestructura de la organización y la falta de un contrato de API diseñado para campos propios del CFDI mexicano como el UUID.

En tercer lugar se encuentra el sistema propuesto en el presente trabajo, el cual combina OCR abierto con preprocesamiento (escala de grises, \textit{deskew}, binarización), clasificación por reglas (factura\textbar orden\_compra\textbar remision\textbar desconocido), persistencia relacional y API REST tipada, con campos fiscales mexicanos (RFC, UUID, totales) como ciudadanos de primera clase del esquema de datos\cite{sat_cfdi_2024}\cite{fastapi_docs}.

Lo que queda sin resolver —y constituye la brecha que aborda el presente trabajo— es un pipeline abierto, de ejecución local y auditable, con campos CFDI nativos, trazabilidad por documento, página, campo y confianza de reconocimiento, y un contrato de API tipado y validado para su integración con sistemas consumidores.


\subsection{Antecedentes de la empresa} \label{sec:antecedentes_empresa} % audit:req=antecedentes_empresa

\begin{guia}[Guía: Antecedentes de la empresa u organismo receptor (Revisión 1)]
\textbf{¿Para qué sirve?} Describe la organización donde realizas la estadía y el proceso que tu proyecto mejora.

\textbf{Debe responder:}
\begin{itemize}[nosep,leftmargin=*]
  \item ¿Quién es el organismo receptor? (giro o sector, ubicación, breve historia, misión y visión).
  \item ¿Cuál es su estructura y en qué área o departamento se realiza tu proyecto? (organigrama opcional).
  \item ¿Cómo se realiza hoy el proceso que vas a mejorar? ¿con qué personas, herramientas o sistemas?
  \item ¿Quién es tu asesor empresarial y cuál es su función?
\end{itemize}

\textbf{Extensión mínima:} 100 palabras.

\textbf{Cuidado:} pide autorización a la empresa para publicar información; no incluyas datos confidenciales ni
datos personales.

\textbf{Errores frecuentes:} copiar la página web de la empresa; no describir el proceso actual (es lo más importante).

\textbf{Cómo se revisa:} que se entienda el contexto operativo real del proyecto.
\end{guia}

% >>> ESCRIBE AQUÍ los antecedentes de la empresa <<<
El Centro de Innovación Industrial en Inteligencia Artificial, CII.IA, se ubica en Monterrey, Nuevo León, México. Según su información pública, su actividad relaciona la inteligencia artificial con necesidades empresariales mediante consultoría, capacitación y prototipado de soluciones (CII.IA, 2026a)\cite{ciiia_2026a}. La organización comunica un enfoque orientado a pasar de la identificación de una necesidad a la construcción, validación y escalamiento de soluciones de inteligencia artificial, además de ofrecer entornos controlados para probar ideas y construir productos mínimos viables (CII.IA, 2026b)\cite{ciiia_2026b}.

En dicho contexto se desarrolló la estadía profesional del autor, cuyo propósito fue construir un prototipo funcional de digitalización inteligente de facturas y documentos de compra. El prototipo corresponde al módulo de OCR, extracción estructurada y backend descrito en la presente tesina, y se concibe como un producto mínimo viable susceptible de validación y escalamiento posterior dentro de los entornos controlados que la organización ofrece\cite{ciiia_2026b}.

El proceso que el prototipo busca mejorar es la captura semimanual de documentos de compra: los documentos se reciben por correo electrónico o en físico, se revisan visualmente y se capturan en hojas de cálculo o sistemas contables mediante transcripción. Este esquema, si bien requiere baja inversión inicial, presenta costos ocultos en tiempos de captura, errores de digitación y dificultad para rastrear el estado de cada documento\cite{manual_entry_errors_2023}.


\subsection{Antecedentes teóricos} \label{sec:antecedentes_teoricos} % audit:req=antecedentes_teoricos

\begin{guia}[Guía: Antecedentes teóricos (Revisión 1)]
\textbf{¿Para qué sirve?} Presenta, de manera panorámica, los conceptos y tecnologías base de tu proyecto. El desarrollo
completo va en el \emph{Marco teórico} (Revisión 2); aquí basta con lo esencial para entender el planteamiento.

\textbf{Debe responder:} ¿qué conceptos, teorías o tecnologías necesita conocer el lector? ¿por qué son relevantes para
tu proyecto? ¿qué fuentes los respaldan?

\textbf{Extensión mínima:} 150 palabras (sin contar la tabla comparativa).

\textbf{Errores frecuentes:} definiciones sueltas sin relación con el proyecto; fuentes sin citar.

\textbf{Cómo se revisa:} pertinencia y suficiencia de los conceptos, y calidad de las citas.
\end{guia}

% >>> ESCRIBE AQUÍ los antecedentes teóricos <<<
\paragraph{Reconocimiento óptico de caracteres (OCR).} El reconocimiento óptico de caracteres (OCR, por sus siglas en inglés) es el proceso mediante el cual una computadora convierte imágenes de texto —ya sean escaneadas, fotografiadas o renderizadas— en caracteres codificados procesables y editables\cite{tesseract_ocr}.

Los motores modernos de OCR operan mediante una secuencia de etapas: preprocesamiento de la imagen, segmentación de líneas y palabras, reconocimiento de caracteres mediante modelos entrenados y asignación de un nivel de confianza a cada decisión. Tesseract, uno de los motores abiertos más utilizados, permite configurar el idioma de reconocimiento y el modo de segmentación de página (PSM, por sus siglas en inglés), así como obtener la confianza de cada palabra reconocida, información que resulta fundamental para dirigir la revisión humana hacia los campos dudosos\cite{tesseract_ocr}.

\paragraph{Extracción estructurada de información.} La extracción estructurada es el proceso de identificar, dentro de un texto no estructurado, un conjunto de campos con significado definido (por ejemplo, folio, fecha, emisor, conceptos, totales) y representarlos conforme a un esquema tipado\cite{structured_extraction_2023}.

En el caso de los documentos de compra, la extracción estructurada enfrenta desafíos particulares: los formatos varían entre emisores, los conceptos se presentan como tablas de renglones múltiples y los montos admiten distintos formatos numéricos. Por ello, los sistemas de extracción combinan expresiones regulares, reglas heurísticas y, en fases avanzadas, modelos de lenguaje, además de reportar métricas de exactitud por campo para demostrar su desempeño\cite{structured_extraction_2023}.

\paragraph{Interfaces de programación de aplicaciones (API REST).} Una interfaz de programación de aplicaciones (API, por sus siglas en inglés) es un contrato mediante el cual un sistema expone sus funcionalidades a otros programas de forma controlada. El estilo arquitectónico REST (transferencia de estado representacional) organiza dicho contrato alrededor de recursos identificados por rutas y operaciones estándar del protocolo HTTP\cite{fielding_rest_2000}.

En el presente proyecto, la API REST constituye el punto de acceso al pipeline de digitalización: a través de ella se cargan documentos, se consultan resultados y se recuperan los campos extraídos. El uso de esquemas de validación tipados (Pydantic v2) garantiza que los datos intercambiados cumplan el contrato antes de persistirse en la base de datos, lo cual reduce errores de integración con sistemas consumidores como el frontend del equipo\cite{fastapi_docs}.


\subsubsection{Tabla comparativa de antecedentes} \label{subsec:tabla_comparativa_antecedentes} % audit:req=tabla_comparativa

\begin{guia}[Guía: Tabla comparativa de antecedentes (Revisión 1)]
\textbf{¿Para qué sirve?} Compara de forma objetiva los trabajos o sistemas existentes con lo que propones.

\textbf{Debe incluir:}
\begin{itemize}[nosep,leftmargin=*]
  \item Una fila por trabajo o sistema (mínimo 3), cada uno con su cita \verb|\cite{clave}|.
  \item Columnas con \emph{criterios} relevantes para tu proyecto (funcionalidades, tecnología, costo, precisión,
        escalabilidad, limitaciones\ldots).
  \item Una fila final con \textbf{tu propuesta}, para que se vea qué aporta de distinto.
  \item Un pie de tabla (\verb|\caption|), una etiqueta (\verb|\label{tab:...}|) y una referencia en el texto
        (\verb|Tabla~\ref{tab:...}|), además de un párrafo que \emph{interprete} la tabla.
\end{itemize}
Usa \verb|booktabs| (\verb|\toprule|, \verb|\midrule|, \verb|\bottomrule|) y evita líneas verticales. Hay un ejemplo
completo en \texttt{ejemplos/ejemplos\_latex.tex}.

\textbf{Errores frecuentes:} tabla sin analizar en el texto; criterios que no ayudan a decidir; no incluir la propuesta.

\textbf{Cómo se revisa:} se verifica automáticamente que exista una tabla; el director valora que los criterios sean
significativos.
\end{guia}

\begin{table}[H]
  \centering
  \small
  \caption{Comparación de enfoques de digitalización frente al sistema propuesto.}
  \label{tab:comparativa_antecedentes}
  \begin{tabular}{@{}p{2.5cm}p{3.0cm}p{3.0cm}p{3.0cm}p{3.0cm}@{}}
    \toprule
    \textbf{Solución} & \textbf{Técnica OCR/extracción} & \textbf{Validación fiscal CFDI} & \textbf{Integración API} & \textbf{Limitación frente a lo propuesto} \\
    \midrule
    Tesseract + reglas clásico &
    OCR Tesseract (spa, PSM 6) con preprocesamiento OpenCV y extracción por expresiones regulares &
    Sin validación fiscal mexicana; campos genéricos (folio, fecha, total) &
    Sin contrato de API definido; scripts o salidas a archivo &
    Reglas frágiles ante formatos no previstos; sin persistencia relacional ni trazabilidad por palabra \\
    \midrule
    Servicio en la nube (Google Cloud Document AI) &
    OCR neuronal propietario y extracción pre-entrenada de facturas y tablas &
    Genérica e internacional; sin campos CFDI (UUID, RFC) como nativos &
    API propietaria por documento con costo asociado &
    Dependencia de proveedor externo y costo por documento; documentos fiscales fuera de la infraestructura local; sin esquema CFDI dedicado \\
    \midrule
    \textbf{Propuesta (este trabajo)} &
    Tesseract (spa, PSM 6, \textit{deskew}) con confianza por palabra + clasificador por reglas + extracción estructurada &
    Campos CFDI nativos: RFC, UUID, folio, conceptos y totales con tolerancia de $\pm 1.00$ &
    API REST propia con FastAPI + Pydantic v2 y persistencia en PostgreSQL; contrato abierto y auditable &
    Ventaja propuesta: ejecución local open-source sin costo por documento, esquema fiscal mexicano y trazabilidad completa (documento, página, campo, confianza) \\
    \bottomrule
  \end{tabular}
\end{table}

La Tabla~\ref{tab:comparativa_antecedentes} muestra que la vía clásica abierta carece de validación fiscal mexicana y de contrato de integración, mientras que la vía de nube resuelve la exactitud a costa de dependencia externa, costo por documento y salida de los documentos fiscales de la infraestructura local. La propuesta se distingue por combinar ejecución local sin costo por documento con campos CFDI nativos, persistencia relacional auditable y un contrato de API tipado, que es precisamente la brecha detectada en los antecedentes del problema.
